\begin{abstract}
Security teams often use small AI text classifiers to sort incoming emails and
text messages. The classifier decides whether a message is \emph{malicious}
(harmful) or \emph{benign} (safe) before a human reviews it. In this paper,
we test what happens when an attacker makes small text edits, called
perturbations, to fool the classifier. We show that evasion rates can only be
measured fairly under a \emph{conditional} test rule: researchers must only
attack messages that the model originally caught, re-check every original
message during testing, and limit the attacker to a realistic edit budget per
sample.

We test this method on three open, lightweight language models (1.7 billion to
3.8 billion parameters) without extra training (zero-shot) on phishing emails
and SMS spam (known as smishing). We test 2{,}079 caught messages using 8{,}047
model queries under a 15\% total edit budget. At first glance, the models look
quite resistant to simple edits. Small character changes (like typos or
look-alike letters) flip only 0.0--5.0\% of caught messages. Simple word swaps
flip only 1.0--5.1\% while changing only 6.2--9.2\% of words, about half the
share of units the character attack spends.

However, a more targeted search changes the picture. When we rank word
replacements by their actual impact on the model's decision (a saliency-guided
word attack), the success rate rises to 4.6\%, compared to 2.7\% for random word
swaps. This guided attack achieves 46 flips across 1{,}001 attempts, averaging 23
model queries per message. In direct paired comparisons on the exact same
messages, the guided attack flips 33 messages that random edits miss, while
random edits flip only 7 that guided edits miss (exact McNemar test
$p=4\times10^{-5}$). A similar but smaller gap appears between word and
character edits (26 versus 10, $p=0.011$). This shows that vulnerability is
concentrated in a few key words rather than spread evenly across the vocabulary.
Evasion rates depend heavily on the quality of the attacker's search strategy.

Two other key findings help interpret these numbers. First, text edits do not
transfer well across different models. When we test saved edits against other
models, only 0.9\% of caught messages flip off the diagonal. Second,
continuous probability scores (the model's internal confidence) do not always
match the final text decision. The two channels agree on 941 of 1{,}001 rows,
matching perfectly for one model and diverging on the other two. We report final
text decisions throughout and treat this gap as an important measurement
finding. For defences, cleaning look-alike characters stops all character-level
evasions and barely dents the guided attack (7.5\% to 6.9\%). Adjusting decision
cutoffs (re-thresholding) to maintain a 5\% false-alarm budget can actually
increase evasions for models with uncalibrated probability scores.

Our main practical takeaway is that under small edit budgets, neither overall
accuracy nor the distance from the decision threshold orders how often a model
flips: the model with the most flips is not the least accurate, and its
probability channel saturates, so no margin can be read on it. Robustness
therefore has to be measured per model rather than inferred from accuracy or
confidence. Furthermore, attack success rates must
always be reported with their exact sample count (denominator), a statistical
confidence interval, and the specific search method used. A result of 0 out of
188 is an upper limit of 2.0\%, not proof of absolute security. We release all
frozen datasets, attack code, and individual sample records to support
reproducible security evaluations.
\end{abstract}